% Preámbulo compartido para todas las clases de Beamer.
% Cada clase lo incluye con % Preámbulo compartido para todas las clases de Beamer.
% Cada clase lo incluye con % Preámbulo compartido para todas las clases de Beamer.
% Cada clase lo incluye con \input{../../latex/preamble.tex}

\documentclass[aspectratio=169,11pt]{beamer}

\usetheme{Madrid}
\setbeamertemplate{navigation symbols}{}
\setbeamertemplate{footline}[frame number]

\usepackage[utf8]{inputenc}
\usepackage[spanish]{babel}
\usepackage{amsmath,amssymb}
\usepackage{graphicx}
\usepackage{booktabs}
\usepackage{listings}
\usepackage{xcolor}

\lstset{
  basicstyle=\ttfamily\small,
  keywordstyle=\color{blue},
  commentstyle=\color{gray},
  stringstyle=\color{teal},
  showstringspaces=false,
  breaklines=true,
  frame=single
}

\newcommand{\definicion}[1]{%
  \begin{block}{Definición}
    #1
  \end{block}
}

\newcommand{\ejemplo}[1]{%
  \begin{exampleblock}{Ejemplo}
    #1
  \end{exampleblock}
}

\newcommand{\importante}[1]{%
  \begin{alertblock}{Importante}
    #1
  \end{alertblock}
}

\author{Agustín Gurvich}
\institute{Introducción a Machine Learning}


\documentclass[aspectratio=169,11pt]{beamer}

\usetheme{Madrid}
\setbeamertemplate{navigation symbols}{}
\setbeamertemplate{footline}[frame number]

\usepackage[utf8]{inputenc}
\usepackage[spanish]{babel}
\usepackage{amsmath,amssymb}
\usepackage{graphicx}
\usepackage{booktabs}
\usepackage{listings}
\usepackage{xcolor}

\lstset{
  basicstyle=\ttfamily\small,
  keywordstyle=\color{blue},
  commentstyle=\color{gray},
  stringstyle=\color{teal},
  showstringspaces=false,
  breaklines=true,
  frame=single
}

\newcommand{\definicion}[1]{%
  \begin{block}{Definición}
    #1
  \end{block}
}

\newcommand{\ejemplo}[1]{%
  \begin{exampleblock}{Ejemplo}
    #1
  \end{exampleblock}
}

\newcommand{\importante}[1]{%
  \begin{alertblock}{Importante}
    #1
  \end{alertblock}
}

\author{Agustín Gurvich}
\institute{Introducción a Machine Learning}


\documentclass[aspectratio=169,11pt]{beamer}

\usetheme{Madrid}
\setbeamertemplate{navigation symbols}{}
\setbeamertemplate{footline}[frame number]

\usepackage[utf8]{inputenc}
\usepackage[spanish]{babel}
\usepackage{amsmath,amssymb}
\usepackage{graphicx}
\usepackage{booktabs}
\usepackage{listings}
\usepackage{xcolor}

\lstset{
  basicstyle=\ttfamily\small,
  keywordstyle=\color{blue},
  commentstyle=\color{gray},
  stringstyle=\color{teal},
  showstringspaces=false,
  breaklines=true,
  frame=single
}

\newcommand{\definicion}[1]{%
  \begin{block}{Definición}
    #1
  \end{block}
}

\newcommand{\ejemplo}[1]{%
  \begin{exampleblock}{Ejemplo}
    #1
  \end{exampleblock}
}

\newcommand{\importante}[1]{%
  \begin{alertblock}{Importante}
    #1
  \end{alertblock}
}

\author{Agustín Gurvich}
\institute{Introducción a Machine Learning}


\title{Sesión 3: Regresión Lineal y Logística}
\subtitle{De predecir números a clasificar categorías}
\date{}

\begin{document}

\frame{\titlepage}

\part{Regresión lineal}
\frame{\partpage}

\section{¿Qué es regresión lineal?}

\begin{frame}{¿Qué es regresión lineal?}
  \definicion{
    Un modelo que predice un valor \textbf{numérico continuo} a partir de
    una combinación lineal de las variables de entrada.
  }
  \ejemplo{
    Precio de una casa a partir de sus m², cantidad de ambientes y
    antigüedad. El precio puede ser cualquier número, no una categoría.
  }
  Es el primer modelo que vamos a entrenar a full.
\end{frame}

\section{Regresión lineal simple}

\begin{frame}{La recta que mejor ajusta}
  Con una sola variable de entrada $x$, buscamos la recta
  $$\hat{y} = w x + b$$
  \begin{itemize}
    \item $w$: pendiente, cuánto cambia $\hat{y}$ por cada unidad de $x$
    \item $b$: ordenada al origen, valor de $\hat{y}$ cuando $x = 0$
  \end{itemize}
  ``Entrenar'' el modelo es eso: encontrar los valores de $w$ y
  $b$ que mejor ajustan a los datos.
\end{frame}

\begin{frame}{¿Qué significa ``mejor ajusta''?}
  \definicion{
    El \textbf{error cuadrático medio} (MSE) mide qué tan lejos están las
    predicciones del modelo de los valores reales.
    $$\text{MSE} = \frac{1}{n}\sum_{i=1}^{n} (y_i - \hat{y}_i)^2$$
  }
  Elevamos al cuadrado para que errores positivos y negativos no se
  cancelen entre sí, y para penalizar más los errores grandes. Entrenar es
  buscar el $w$ y el $b$ que minimizan el MSE.
\end{frame}

\section{Descenso de gradiente}

\begin{frame}{¿Cómo encontramos el mínimo?}
  \begin{enumerate}
    \item Empezar con $w$ y $b$ al azar (predicciones malas)
    \item Medir el error
    \item Mover $w$ y $b$ un poquito en la dirección que lo reduce
    \item Repetir hasta que el error deje de bajar
  \end{enumerate}
  \pause
  \ejemplo{
    ¿Y cómo salimos de este pozo? Pues cavando.
  }
\end{frame}

\begin{frame}{Learning rate}
  \begin{itemize}
    \item Controla el tamaño de cada paso del descenso de gradiente
    \item Muy chico: converge, pero tarda muchísimas iteraciones
    \item Muy grande: puede pasarse del mínimo y no converger nunca
  \end{itemize}
  \pause
  \importante{
    Nada de esto se programa a mano: scikit-learn ya lo resuelve,
    y de forma más eficiente que cualquier loop que escribamos nosotros.
  }
\end{frame}

\section{Regresión lineal múltiple}

\begin{frame}{Más de una variable}
  En un problema real casi siempre hay varias features. La idea es la
  misma, con un coeficiente por cada una:
  $$\hat{y} = w_1 x_1 + w_2 x_2 + \dots + w_n x_n + b$$
  \ejemplo{
    Precio $=$ $w_1 \cdot$ m² $+ w_2 \cdot$ ambientes $+ w_3 \cdot$
    antigüedad $+ b$
  }
\end{frame}

\begin{frame}{Interpretando los coeficientes}
  \begin{itemize}
    \item Coeficiente grande y positivo: esa variable empuja fuerte la
      predicción para arriba
    \item Coeficiente negativo: la empuja para abajo
    \item Coeficiente cercano a 0: esa variable casi no influye
  \end{itemize}
  \pause
  Ojo: si las variables están en escalas muy distintas, comparar
  coeficientes a simple vista engaña. Por eso escalamos antes de entrenar!
\end{frame}

\section{Evaluando el ajuste}

\begin{frame}{$R^2$: coeficiente de determinación}
  \definicion{
    Mide qué proporción de la variabilidad de $y$ explica el modelo. Va de
    $-\infty$ a $1$.
  }
  \begin{itemize}
    \item $R^2 = 1$: predicciones perfectas
    \item $R^2 = 0$: el modelo predice tan bien como usar siempre el
      promedio de $y$
    \item $R^2 < 0$: el modelo es peor que predecir el promedio
  \end{itemize}
\end{frame}

\section{Regresión lineal en código}

\begin{frame}[fragile]{Regresión lineal con scikit-learn}
  \begin{lstlisting}[language=Python]
from sklearn.linear_model import LinearRegression

modelo = LinearRegression()
modelo.fit(X_train, y_train)

modelo.coef_        # los w_i
modelo.intercept_   # el b

modelo.predict(X_test)
modelo.score(X_test, y_test)   # R^2
  \end{lstlisting}
\end{frame}

\part{Regresión logística}
\frame{\partpage}

\section{De regresión a clasificación}

\begin{frame}{El problema: clasificar, no predecir un número}
  \ejemplo{
    Si un mail es spam o no, si un tumor es maligno o benigno, si un
    cliente deja de pagar o no. La salida es una \textbf{categoría}, no un
    número continuo.
  }
  ¿Por qué no usar regresión lineal directamente, prediciendo 0 o 1?
  \pause
  \importante{
    Porque puede devolver cualquier número: $-3$, $1.7$, $250$. Eso no se
    interpreta como ``sí'' o ``no'', y un solo outlier corre toda la
    recta.
  }
\end{frame}

\section{La función sigmoide}

\begin{frame}{La función sigmoide}
  \definicion{
    Aplasta cualquier número real al rango $(0, 1)$.
    $$\sigma(z) = \frac{1}{1 + e^{-z}}$$
  }
  \begin{itemize}
    \item $z \to +\infty \Rightarrow \sigma(z) \to 1$
    \item $z \to -\infty \Rightarrow \sigma(z) \to 0$
    \item $z = 0 \Rightarrow \sigma(z) = 0.5$
  \end{itemize}
  Forma de ``S'': un número cualquiera entra, algo interpretable como
  probabilidad sale.
\end{frame}

\section{Regresión logística: el modelo}

\begin{frame}{Combinando todo}
  Regresión lineal, pasada por la sigmoide:
  $$\hat{p} = \sigma(w_1 x_1 + \dots + w_n x_n + b)$$
  \definicion{
    $\hat{p}$ se interpreta como la \textbf{probabilidad} de que la
    observación sea de la clase positiva (por ejemplo, ``es spam'').
  }
\end{frame}

\begin{frame}{De probabilidad a clase: el umbral}
  \begin{itemize}
    \item $\hat{p} \geq 0.5$: predecimos clase 1
    \item $\hat{p} < 0.5$: predecimos clase 0
  \end{itemize}
  \pause
  \importante{
    $0.5$ es el umbral por defecto, no una ley física. En diagnóstico
    médico, por ejemplo, conviene bajarlo para no dejar pasar casos
    positivos, aunque eso traiga más falsos positivos.
  }
\end{frame}

\section{Entrenamiento: log loss}

\begin{frame}{¿Cómo se entrena?}
  \definicion{
    La función de costo acá es la \textbf{log loss}: penaliza fuerte
    cuando el modelo está muy seguro y se equivoca.
  }
  \ejemplo{
    Predecir $\hat{p} = 0.99$ para algo que en realidad es clase 0 sale
    mucho más caro que predecir $\hat{p} = 0.6$ y fallar igual.
  }
\end{frame}

\section{Más de dos clases}

\begin{frame}{¿Y si hay más de dos categorías?}
  \begin{itemize}
    \item \textbf{One-vs-rest}: un clasificador binario por clase (``es
      esta'' vs. ``no lo es''), se queda con el de mayor probabilidad
    \item \textbf{Softmax}: generaliza la sigmoide a varias clases a la
      vez, con probabilidades que suman 1
  \end{itemize}
  scikit-learn elige la estrategia solo, según el problema.
\end{frame}

\section{Regresión logística en código}

\begin{frame}[fragile]{Regresión logística con scikit-learn}
  \begin{lstlisting}[language=Python]
from sklearn.linear_model import LogisticRegression

modelo = LogisticRegression()
modelo.fit(X_train, y_train)

modelo.predict(X_test)          # clases predichas
modelo.predict_proba(X_test)    # probabilidad por clase

modelo.score(X_test, y_test)    # accuracy
  \end{lstlisting}
  \pause
  \texttt{predict\_proba} es la gran diferencia frente a regresión lineal:
  no solo da una clase, da qué tan segura está.
\end{frame}

\begin{frame}{Resumen}
  \begin{itemize}
    \item Regresión lineal: ajusta una recta (o hiperplano) minimizando el
      error cuadrático medio
    \item Regresión logística: la misma idea, pasada por una sigmoide, para
      obtener una probabilidad
    \item Ambas se entrenan con descenso de gradiente y comparten la
      interfaz \texttt{fit} / \texttt{predict} de scikit-learn
  \end{itemize}
  \vspace{1em}
  \centering
  \Large ¿Preguntas?
\end{frame}

\end{document}
